\section{Blow up algebras}
\label{section-blow-up}

\noindent
In this section we make some elementary observations about blowing up.

\begin{definition}
\label{definition-blow-up}
Let $R$ be a ring.
Let $I \subset R$ be an ideal.
\begin{enumerate}
\item The {\it blowup algebra}, or the {\it Rees algebra}, associated to
the pair $(R, I)$ is the graded $R$-algebra
$$
\text{Bl}_I(R) =
\bigoplus\nolimits_{n \geq 0} I^n =
R \oplus I \oplus I^2 \oplus \ldots
$$
where the summand $I^n$ is placed in degree $n$.
\item Let $a \in I$ be an element. Denote $a^{(1)}$ the element $a$
seen as an element of degree $1$ in the Rees algebra. Then the
{\it affine blowup algebra} $R[\frac{I}{a}]$ is the algebra
$(\text{Bl}_I(R))_{(a^{(1)})}$ constructed in Section \ref{section-proj}.
\end{enumerate}
\end{definition}

\noindent
In other words, an element of $R[\frac{I}{a}]$ is represented by
an expression of the form $x/a^n$ with $x \in I^n$. Two representatives
$x/a^n$ and $y/a^m$ define the same element if and only if
$a^k(a^mx - a^ny) = 0$ for some $k \geq 0$.

\begin{lemma}
\label{lemma-affine-blowup}
Let $R$ be a ring, $I \subset R$ an ideal, and $a \in I$.
Let $R' = R[\frac{I}{a}]$ be the affine blowup algebra. Then
\begin{enumerate}
\item the image of $a$ in $R'$ is a nonzerodivisor,
\item $IR' = aR'$, and
\item $(R')_a = R_a$.
\end{enumerate}
\end{lemma}

\begin{proof}
Immediate from the description of $R[\frac{I}{a}]$ above.
\end{proof}

\begin{lemma}
\label{lemma-blowup-base-change}
Let $R \to S$ be a ring map. Let $I \subset R$ be an ideal
and $a \in I$. Set $J = IS$ and let $b \in J$ be the image of $a$.
Then $S[\frac{J}{b}]$ is the quotient of $S \otimes_R R[\frac{I}{a}]$
by the ideal of elements annihilated by some power of $b$.
\end{lemma}

\begin{proof}
Let $S'$ be the quotient of $S \otimes_R R[\frac{I}{a}]$ by its
$b$-power torsion elements. The ring map
$$
S \otimes_R R[\textstyle{\frac{I}{a}}]
\longrightarrow
S[\textstyle{\frac{J}{b}}]
$$
is surjective and annihilates $a$-power torsion as $b$ is a nonzerodivisor
in $S[\frac{J}{b}]$. Hence we obtain a surjective map $S' \to S[\frac{J}{b}]$.
To see that the kernel is trivial, we construct an inverse map. Namely, let
$z = y/b^n$ be an element of $S[\frac{J}{b}]$, i.e., $y \in J^n$.
Write $y = \sum x_is_i$ with $x_i \in I^n$ and $s_i \in S$.
We map $z$ to the class of $\sum s_i \otimes x_i/a^n$ in
$S'$. This is well defined because an element of the kernel of the map
$S \otimes_R I^n \to J^n$ is annihilated by $a^n$, hence maps to zero in $S'$.
\end{proof}

\begin{example}
\label{example-rees-algebra-polynomial}
Let $R$ be a ring. Let $P = R[t_1, \ldots, t_n]$ be the polynomial algebra.
Let $I = (t_1, \ldots, t_n) \subset P$. With notation as in
Definition \ref{definition-blow-up} there is an isomorphism
$$
P[T_1, \ldots, T_n]/(t_iT_j - t_jT_i) \longrightarrow \text{Bl}_I(P)
$$
sending $T_i$ to $t_i^{(1)}$. We leave it to the reader to show that
this map is well defined. Since $I$ is generated by $t_1, \ldots, t_n$
we see that our map is surjective. To see that our map is injective
one has to show: for each $e \geq 1$ the $P$-module $I^e$ is generated
by the monomials $t^E = t_1^{e_1} \ldots x_n^{e_n}$ for multiindices
$E = (e_1, \ldots, e_n)$ of degree $|E| = e$ subject only to the
relations $t_i t^E = t_j t^{E'}$ when $|E| = |E'| = e$ and
$e_a + \delta_{a i} = e'_a + \delta_{a j},\ a = 1, \ldots, n$
(Kronecker delta). We omit the details.
\end{example}

\begin{example}
\label{example-affine-blowup-algebra-polynomial}
Let $R$ be a ring. Let $P = R[t_1, \ldots, t_n]$ be the polynomial algebra.
Let $I = (t_1, \ldots, t_n) \subset P$. Let $a = t_1$. With notation as in
Definition \ref{definition-blow-up} there is an isomorphism
$$
P[x_2, \ldots, x_n]/(t_1x_2 - t_2, \ldots, t_1x_n - t_n)
\longrightarrow
\textstyle{P[\frac{I}{a}] = P[\frac{I}{t_1}]}
$$
sending $x_i$ to $t_i/t_1$. We leave it to the reader to show that
this map is well defined. Since $I$ is generated by $t_1, \ldots, t_n$
we see that our map is surjective. To see that our map is injective,
the reader can argue that the source and target of our map are
$t_1$-torsion free and that the map is an isomorphism after inverting
$t_1$, see Lemma \ref{lemma-affine-blowup}. Alternatively, the reader
can use the description of the Rees algebra in
Example \ref{example-rees-algebra-polynomial}.
We omit the details.
\end{example}

\begin{lemma}
\label{lemma-affine-blowup-quotient-description}
Let $R$ be a ring. Let $I = (a_1, \ldots, a_n)$ be an ideal of $R$.
Let $a = a_1$. Then there is a surjection
$$
R[x_2, \ldots, x_n]/(a x_2 - a_2, \ldots, a x_n - a_n)
\longrightarrow
\textstyle{R[\frac{I}{a}]}
$$
whose kernel is the $a$-power torsion in the source.
\end{lemma}

\begin{proof}
Consider the ring map $P = \mathbf{Z}[t_1, \ldots, t_n] \to R$
sending $t_i$ to $a_i$. Set $J = (t_1, \ldots, t_n)$.
By Example \ref{example-affine-blowup-algebra-polynomial} we have
$P[\frac{J}{t_1}] =
P[x_2, \ldots, x_n]/(t_1 x_2 - t_2, \ldots, t_1 x_n - t_n)$.
Apply Lemma \ref{lemma-blowup-base-change} to the map $P \to A$ to conclude.
\end{proof}

\begin{lemma}
\label{lemma-blowup-in-principal}
Let $R$ be a ring, $I \subset R$ an ideal, and $a \in I$.
Set $R' = R[\frac{I}{a}]$. If $f \in R$ is such that $V(f) = V(I)$,
then $f$ maps to a nonzerodivisor in $R'$ and $R'_f = R'_a = R_a$.
\end{lemma}

\begin{proof}
We will use the results of Lemma \ref{lemma-affine-blowup}
without further mention.
The assumption $V(f) = V(I)$ implies $V(fR') = V(IR') = V(aR')$.
Hence $a^n = fb$ and $f^m = ac$ for some $b, c \in R'$.
The lemma follows.
\end{proof}

\begin{lemma}
\label{lemma-blowup-add-principal}
Let $R$ be a ring, $I \subset R$ an ideal, $a \in I$, and $f \in R$.
Set $R' = R[\frac{I}{a}]$ and $R'' = R[\frac{fI}{fa}]$. Then
there is a surjective $R$-algebra map $R' \to R''$ whose kernel
is the set of $f$-power torsion elements of $R'$.
\end{lemma}

\begin{proof}
The map is given by sending $x/a^n$ for $x \in I^n$ to $f^nx/(fa)^n$.
It is straightforward to check this map is well defined and surjective.
Since $af$ is a nonzero divisor in $R''$
(Lemma \ref{lemma-affine-blowup}) we see that the set of $f$-power
torsion elements are mapped to zero. Conversely, if $x \in R'$
and $f^n x \not = 0$ for all $n > 0$, then $(af)^n x \not = 0$
for all $n$ as $a$ is a nonzero divisor in $R'$. It follows
that the image of $x$ in $R''$ is not zero by the description of
$R''$ following Definition \ref{definition-blow-up}.
\end{proof}

\begin{lemma}
\label{lemma-blowup-reduced}
\begin{slogan}
Being reduced is invariant under blowup
\end{slogan}
If $R$ is reduced then every (affine) blowup algebra of $R$ is reduced.
\end{lemma}

\begin{proof}
Let $I \subset R$ be an ideal and $a \in I$. Suppose $x/a^n$ with
$x \in I^n$ is a nilpotent element of $R[\frac{I}{a}]$. Then
$(x/a^n)^m = 0$. Hence $a^N x^m = 0$ in $R$ for some $N \geq 0$.
After increasing $N$ if necessary we may assume $N = me$ for some
$e \geq 0$. Then $(a^e x)^m = 0$ and since $R$ is reduced we find
$a^e x = 0$. This means that $x/a^n = 0$ in $R[\frac{I}{a}]$.
\end{proof}

\begin{lemma}
\label{lemma-blowup-domain}
Let $R$ be a domain, $I \subset R$ an ideal, and $a \in I$ a nonzero
element. Then the affine blowup algebra $R[\frac{I}{a}]$ is a domain.
\end{lemma}

\begin{proof}
Suppose $x/a^n$, $y/a^m$ with $x \in I^n$, $y \in I^m$
are elements of $R[\frac{I}{a}]$ whose product is zero.
Then $a^N x y = 0$ in $R$. Since $R$ is a domain we conclude
that either $x = 0$ or $y = 0$.
\end{proof}

\begin{lemma}
\label{lemma-blowup-dominant}
Let $R$ be a ring. Let $I \subset R$ be an ideal. Let $a \in I$.
If $a$ is not contained in any minimal prime of $R$, then
$\Spec(R[\frac{I}{a}]) \to \Spec(R)$ has dense image.
\end{lemma}

\begin{proof}
If $a^k x = 0$ for $x \in R$, then $x$ is contained in all the
minimal primes of $R$ and hence nilpotent, see
Lemma \ref{lemma-Zariski-topology}.
Thus the kernel of $R \to R[\frac{I}{a}]$ consists of nilpotent
elements. Hence the result follows from
Lemma \ref{lemma-image-dense-generic-points}.
\end{proof}

\begin{lemma}
\label{lemma-valuation-ring-colimit-affine-blowups}
Let $(R, \mathfrak m)$ be a local domain with fraction field $K$.
Let $R \subset A \subset K$ be a valuation ring which dominates $R$.
Then
$$
A = \colim R[\textstyle{\frac{I}{a}}]
$$
is a directed colimit of affine blowups $R \to R[\frac{I}{a}]$ with
the following properties
\begin{enumerate}
\item $a \in I \subset \mathfrak m$,
\item $I$ is finitely generated, and
\item the fibre ring of $R \to R[\frac{I}{a}]$ at $\mathfrak m$
is not zero.
\end{enumerate}
\end{lemma}

\begin{proof}
Any blowup algebra $R[\frac{I}{a}]$ is a domain contained in $K$ see
Lemma \ref{lemma-blowup-domain}. The lemma simply says that $A$ is the
directed union of the ones where $a \in I$ have properties (1), (2), (3).
If $R[\frac{I}{a}] \subset A$ and $R[\frac{J}{b}] \subset A$, then
we have
$$
R[\textstyle{\frac{I}{a}}] \cup R[\textstyle{\frac{J}{b}}]  \subset
R[\textstyle{\frac{IJ}{ab}}] \subset A
$$
The first inclusion because $x/a^n = b^nx/(ab)^n$ and the second one
because if $z \in (IJ)^n$, then $z = \sum x_iy_i$ with $x_i \in I^n$
and $y_i \in J^n$ and hence $z/(ab)^n = \sum (x_i/a^n)(y_i/b^n)$
is contained in $A$.

\medskip\noindent
Consider a finite subset $E \subset A$. Say $E = \{e_1, \ldots, e_n\}$.
Choose a nonzero $a \in R$ such that we can write $e_i = f_i/a$ for
all $i = 1, \ldots, n$. Set $I = (f_1, \ldots, f_n, a)$.
We claim that $R[\frac{I}{a}] \subset A$. This is clear as an element
of $R[\frac{I}{a}]$ can be represented as a polynomial in the elements
$e_i$. The lemma follows immediately from this observation.
\end{proof}










